\section[Conclusão]{Conclusão}

O projeto SalesPilot propôs um desafio concreto e relevante: construir uma plataforma de gestão capaz de traduzir os dados gerados por uma ferramenta de inteligência artificial em visualizações claras e acionáveis para três perfis de usuários distintos, administradores, gestores e vendedores. Participar do desenvolvimento dessa plataforma como \ac{ages} III, com responsabilidade integral pela arquitetura do frontend, foi uma experiência que testou minhas capacidades técnicas, minha capacidade de liderança e também minha resiliência diante de um ambiente de trabalho dinâmico e nem sempre previsível.

Do ponto de vista técnico, avalio minha atuação de forma positiva. Desde a Sprint 0, assumi a definição da \textit{technology stack}, a concepção da arquitetura baseada em \textit{feature-based architecture} e a estruturação do repositório, e essas decisões se mostraram sólidas ao longo de todo o semestre. A refatoração ampla de pastas que fiz na Sprint 2 organizou bem o código, mas também causou conflitos de \textit{merge} para colegas com tarefas em andamento. Isso me ensinou que mudanças estruturais grandes precisam ser mais bem comunicadas e, quando possível, divididas em partes menores. Já na Sprint 4, criar do zero uma suíte de testes \textit{end-to-end} com Playwright foi uma contribuição da qual me orgulho bastante, porque mostra que me importei com a qualidade do produto como um todo, e não só com as funcionalidades pelas quais eu era diretamente responsável. O que eu acho que poderia ter sido melhor foi integrar as \textit{branches} de forma mais incremental ao longo das sprints, em vez de deixar tudo concentrado no final, o que gerou inconsistências difíceis de identificar e corrigir com o prazo apertado.

Em relação às \textit{soft skills}, o principal crescimento foi aprender a exercer um papel de referência técnica dentro de um time multidisciplinar. Esse papel exige tanto segurança para tomar decisões sozinha quanto clareza para explicá-las a colegas com níveis de experiência bem diferentes do meu. A experiência mais marcante nesse sentido foi a revisão de \textit{pull requests}: quando percebi que estava virando um gargalo para o time na Sprint 3, não hesitei em levantar o problema na retrospectiva, faltavam prazos intermediários para o fechamento das tarefas. Ver essa mudança ser adotada pelo time na Sprint 4, com resultados concretos, confirmou para mim que apontar problemas de processo com assertividade também faz parte do trabalho, não só resolvê-los na prática.

A relação entre as disciplinas que cursei em paralelo e minha experiência na \ac{ages} foi menos direta do que eu esperava, porque atuei só no frontend, e as competências de que precisei vieram mais de pesquisa própria, documentação técnica e experimentação prática no próprio projeto do que do conteúdo das aulas. Ainda assim, a base de disciplinas como arquitetura de software me ajudou a embasar as escolhas de padrões que tornaram o projeto escalável e fácil de manter, mesmo que essa aplicação tenha sido bem mais intuitiva do que algo que eu tirei direto das aulas.

O SalesPilot não foi o projeto mais memorável em que já trabalhei, mas foi, sem dúvida, aquele em que mais cresci como profissional. Outros projetos trouxeram experiências mais fluidas ou equipes mais alinhadas, mas poucos exigiram o nível de responsabilidade técnica e de liderança que assumi aqui. O escopo real do produto, a presença de \textit{stakeholders} genuinamente investidos e a complexidade de coordenar contribuições de múltiplas turmas com diferentes níveis de maturidade tornaram esse projeto um ambiente de aprendizado denso e valioso.

Como principais lições aprendidas, destaco três pontos. Primeiro, a importância de testar o produto de forma integrada com antecedência, e não só perto da entrega, porque problemas descobertos tarde demandam uma disponibilidade de equipe que nem sempre existe. Segundo, a necessidade de integrar as \textit{branches} com mais frequência, para que inconsistências sejam identificadas mais cedo. Terceiro, o valor de estabelecer prazos claros para tarefas intermediárias, o que abre espaço para revisões com calma e ciclos de correção que melhoram a qualidade final do produto. Termino minha experiência como \ac{ages} III certa de que as escolhas de arquitetura bem definidas desde o início contribuíram muito para o bom andamento do projeto, e mais confiante da minha capacidade de atuar como arquiteta de software. Além disso, me sinto realizada por termos feito uma entrega final excelente, sem grandes intercorrências.
